\chapter{Preliminares matemáticos}
\label{cap:preliminares}

Este capítulo reúne los resultados que se usarán después: el modelo probabilístico
de la supervivencia (\cref{sec:prob}), el análisis convexo y no suave
(\cref{sec:nosuave}) y la teoría de convergencia de métodos iterativos
(\cref{sec:convergencia}). La conjunción de los dos primeros, que la literatura
presenta por separado, es lo que permite el análisis del \cref{cap:estructura}.

\section{El modelo probabilístico de la supervivencia}
\label{sec:prob}

\subsection{Tiempo futuro de vida}

\begin{definition}[Tiempo futuro de vida]
Sea $(\Omega,\mathcal{F},\Prob)$ un espacio de probabilidad. Para un individuo de
edad $x$ se define el \emph{tiempo futuro de vida} como la variable aleatoria
\[
  T_x:\Omega\longrightarrow[0,\infty),
\]
donde $T_x(\omega)$ es el tiempo transcurrido desde la edad $x$ hasta el
fallecimiento en el estado $\omega$.
\end{definition}

\begin{definition}[Funciones de supervivencia y de fallecimiento]
Para $t\ge 0$ se definen
\[
  \npx{t}{x} \;=\; S_x(t) \;=\; \Prob(T_x>t),
  \qquad
  \nqx{t}{x} \;=\; 1-\npx{t}{x} \;=\; F_{T_x}(t).
\]
\end{definition}

La función $\npx{t}{x}$ es no creciente, con $\npx{0}{x}=1$ y
$\lim_{t\to\infty}\npx{t}{x}=0$. La propiedad que se usará repetidamente es la
descomposición multiplicativa.

\begin{lemma}[Descomposición de la supervivencia]
\label{lem:descomposicion}
Para $s,t\ge0$,
\[
  \npx{t+s}{x} \;=\; \npx{t}{x}\cdot \npx{s}{x+t}.
\]
\end{lemma}
\begin{proof}
Por definición de probabilidad condicional,
$\Prob(T_x>t+s)=\Prob(T_x>t)\,\Prob(T_x>t+s \mid T_x>t)$, y el segundo factor es
$\npx{s}{x+t}$ cuando ambas probabilidades proceden de la misma tabla por edad.
En un modelo de cohorte la identidad condicional se conserva, incorporando
también la fecha calendario al segundo factor \citep{pitacco2009}.
\end{proof}

\begin{corollary}[Versión discreta]
\label{cor:discreta}
Si $q_{x}$ denota la probabilidad de fallecer entre las edades $x$ y $x+1$,
entonces
\begin{equation}
\label{eq:npx-producto}
  \npx{t}{x} \;=\; \prod_{j=1}^{t}\bigl(1-q_{x+j-1}\bigr),\qquad t\in\N.
\end{equation}
\end{corollary}

La ecuación \cref{eq:npx-producto} es la forma en que efectivamente se calcula la
supervivencia en el modelo implementado, sobre una malla mensual\footnote{La mensualización de una tabla anual exige una hipótesis sobre la distribución de los fallecimientos dentro del año. El modelo implementado adopta la de \emph{fallecimientos uniformemente distribuidos}, que conduce a la expresión $q^{(t)}_{x}=(q_x/12)/(1-s\,q_x)$ con $s$ la fracción de año transcurrida. Es la convención estándar; véase \citet[cap.~3]{dickson2020}.}.

La \emph{fuerza de mortalidad} es $\mu_{x+t}=-\frac{d}{dt}\ln\npx{t}{x}$, de modo que
$\npx{t}{x}=\exp\bigl(-\int_0^t\mu_{x+s}\,ds\bigr)$. Una tabla sigue la \emph{ley de
Gompertz} si $\mu_x=Ae^{bx}$ con $A,b>0$, lo que equivale aproximadamente a que
$\ln q_x$ sea afín en $x$; esa ley permite contrastar las tablas usadas, como se hace
en la \cref{sec:var-demograficas}.

\subsection{Modelos multivida}

Sea $Y=\{y_1,\dots,y_n\}$ un grupo de personas con tiempos futuros de vida
$T_{y_1},\dots,T_{y_n}$. Sea $D_y$ el tiempo de vencimiento del derecho de $y$,
en las mismas unidades que $T_y$, con $D_y=\infty$ para un derecho vitalicio.
Para derechos de duración determinista se define $H_y=\min\{T_y,D_y\}$ y
$p_y^{\mathrm{der}}(t)=\Prob(H_y>t)$, incorporando la convención mensual de
$\DURA_y$ al fijar el extremo de vigencia. Un contrato \emph{de último
sobreviviente con derecho} paga mientras al menos uno conserve el derecho;
su tiempo de terminación es
\[
  T_Y \;=\; \max_{y\in Y} H_y.
\]

\begin{proposition}[Factor de contingencia del amparo básico]
\label{prop:conab}
Si los tiempos $T_y$, $y\in Y$, son independientes y los $D_y$ son
deterministas, entonces
\begin{equation}
\label{eq:conab}
  \CONAB(t)\;=\;\Prob(T_Y>t)\;=\;1-\prod_{y\in Y}\bigl(1-p_y^{\mathrm{der}}(t)\bigr).
\end{equation}
\end{proposition}
\begin{proof}
$\{T_Y>t\}=\bigcup_y\{H_y>t\}$. Como los $D_y$ son deterministas, la
independencia de los $T_y$ implica la de los $H_y$. La probabilidad del
complemento es, por tanto, $\prod_y\Prob(H_y\le t)
=\prod_y(1-p_y^{\mathrm{der}}(t))$.
\end{proof}

En las fórmulas del motor, $\npx{t}{y}$ denota esta supervivencia efectiva
después de truncar por la duración del derecho; coincide con la supervivencia
biológica mientras no opere ese truncamiento.

\begin{remark}[Sobre la hipótesis de independencia]
La independencia entre las vidas de un grupo familiar es una simplificación del
motor: la literatura documenta correlación positiva en la mortalidad conyugal
\citep{frees1996}. Para el objeto de este trabajo la hipótesis es inocua en un
sentido preciso: cualquier modelo de dependencia produce otro $\CONAB(t)$, que no
depende de $M$ y entra en $\VP$ como un peso, de modo que la forma funcional
estudiada en el \cref{cap:estructura} se conserva; las propiedades que exigen
pesos positivos deben verificarse de nuevo (\cref{sec:var-dependencia}).
\end{remark}

\subsection{Valor presente de flujos contingentes}

\begin{definition}
Sea $\{C_t\}_{t\ge1}$ una sucesión de flujos aleatorios y $v_t>0$ el factor de
descuento del periodo $t$. El \emph{valor presente aleatorio} es
$Z=\sum_{t\ge1}v_tC_t$ y el \emph{valor presente actuarial} es $\Esp[Z]$.
\end{definition}

\begin{theorem}[Intercambio de esperanza y suma]
\label{thm:tonelli}
Si $C_t\ge0$ casi seguramente para todo $t$, entonces
\begin{equation}
\label{eq:tonelli}
  \Esp\Bigl[\sum_{t\ge1}v_tC_t\Bigr]\;=\;\sum_{t\ge1}v_t\,\Esp[C_t],
\end{equation}
con ambos miembros posiblemente infinitos. Si además
$\sum_t v_t\Esp|C_t|<\infty$, ambos miembros son finitos y la serie converge
absolutamente.
\end{theorem}
\begin{proof}
La primera afirmación es el teorema de Tonelli aplicado a la medida producto
$\Prob\otimes\nu$, donde $\nu$ es la medida de conteo en $\N$, para el integrando
no negativo $(\omega,t)\mapsto v_tC_t(\omega)$. La segunda es el teorema de Fubini
bajo la hipótesis de integrabilidad \citep[cap.~18]{billingsley1995}.
\end{proof}

En el modelo de este trabajo los flujos son no negativos y la suma es
\emph{finita} --- el horizonte es de \datoNMeses{} meses ---, de modo que
\cref{eq:tonelli} se aplica sin hipótesis de convergencia adicionales.

\begin{corollary}[Valor presente de una renta contingente]
\label{cor:vp-renta}
Si $C_t=M_t\,\ind_{\{T_Y>t\}}$ con $M_t$ determinista, entonces
\[
  \Esp[Z]\;=\;\sum_{t\ge1}v_t\,M_t\,\CONAB(t).
\]
\end{corollary}
\begin{proof}
$\Esp[\ind_{\{T_Y>t\}}]=\Prob(T_Y>t)=\CONAB(t)$ y $M_t$, $v_t$ son constantes;
se aplica \cref{thm:tonelli}.
\end{proof}

\section{Análisis convexo y no suave}
\label{sec:nosuave}

\subsection{Convexidad y operadores que la preservan}

\begin{lemma}[Operaciones que preservan convexidad]
\label{lem:preserva-convexidad}
\begin{enumerate}[label=(\alph*),nosep]
  \item El máximo puntual de una familia de funciones afines es convexo.
  \item Una combinación cónica $\sum_i w_ig_i$ con $w_i\ge0$ de funciones convexas
        es convexa.
  \item El \textbf{mínimo} de dos funciones convexas \emph{no} es convexo en
        general.
\end{enumerate}
\end{lemma}

El apartado (c) es el que produce la pérdida de convexidad del
\cref{thm:no-convexidad}: el operador de recorte
$\med(a,u,b)=\min\bigl(\max(u,a),b\bigr)$ aplica un mínimo sobre una función
convexa.

\begin{definition}[Operador de recorte]
Para $a\le b$,
\[
  \med(a,u,b)\;=\;\min\bigl(\max(u,a),\,b\bigr)
  \;=\;\begin{cases} a, & u<a,\\ u, & a\le u\le b,\\ b, & u>b.\end{cases}
\]
\end{definition}

\begin{lemma}
\label{lem:recorte}
$\med(a,\cdot,b)$ es no decreciente, $1$-Lipschitz, lineal a trozos con quiebres
en $a$ y $b$; es convexa en $(-\infty,b]$ y cóncava en $[a,\infty)$, y no es ni
convexa ni cóncava en $\R$ si $a<b$.
\end{lemma}
\begin{proof}
La monotonía y la constante de Lipschitz son inmediatas del hecho de que la
pendiente vale $0$ o $1$. Para la última afirmación basta observar que la
pendiente pasa de $0$ a $1$ en $a$ (salto positivo, compatible con convexidad) y
de $1$ a $0$ en $b$ (salto negativo, incompatible con convexidad).
\end{proof}

\subsection{Funciones lineales a trozos}

\begin{definition}[Función lineal a trozos]
$g:\R\to\R$ es \emph{lineal a trozos} (o afín a trozos) si existe una partición
finita $-\infty=\beta_0<\beta_1<\cdots<\beta_m<\beta_{m+1}=+\infty$ tal que $g$ es
afín en cada $[\beta_j,\beta_{j+1}]$. Los puntos $\beta_j$ se llaman
\emph{quiebres}.
\end{definition}

\begin{lemma}
\label{lem:pl-cerrado}
La clase de funciones lineales a trozos es cerrada bajo sumas finitas,
multiplicación por escalares, $\max$, $\min$ y composición. El conjunto de
quiebres de una suma está contenido en la unión de los conjuntos de quiebres de
los sumandos.
\end{lemma}

\subsection{Subdiferenciales}

Para $g$ localmente Lipschitz, sea $D_g$ el conjunto (de medida plena, por el
teorema de Rademacher) donde $g$ es diferenciable.

\begin{definition}[$B$-subdiferencial y subdiferencial de Clarke]
\[
  \dB g(u)\;=\;\Bigl\{\, \lim_{k\to\infty} g'(u_k)\;:\; u_k\to u,\ u_k\in D_g \,\Bigr\},
  \qquad
  \dC g(u)\;=\;\operatorname{conv}\bigl(\dB g(u)\bigr).
\]
\end{definition}

\begin{lemma}
\label{lem:subdif-pl}
Si $g$ es lineal a trozos con pendientes $s^-$ y $s^+$ a izquierda y derecha de
$u$, entonces $\dB g(u)=\{s^-,s^+\}$ y $\dC g(u)=[\min(s^-,s^+),\max(s^-,s^+)]$.
En particular $\dB g(u)=\dC g(u)=\{g'(u)\}$ si $u$ no es un quiebre.
\end{lemma}

Nótese que para el algoritmo importa $\dB$ y no $\dC$: el método de Newton
semisuave usa un elemento de $\dB$, que es una \emph{pendiente efectivamente
realizada} por la función, mientras que un elemento interior de $\dC$ puede no
corresponder a ninguna dirección de la gráfica. Esta distinción explica el sesgo
de las diferencias finitas centradas que se cuantifica en la
\cref{prop:sesgo-diferencias}.

\subsection{Semisuavidad}

\begin{definition}[Semisuavidad, \citealp{qi1993}]
$g$ localmente Lipschitz es \emph{semisuave} en $u$ si es direccionalmente
diferenciable en $u$ y
\[
  \sup_{V\in\dC g(u+d)} \bigl| g(u+d)-g(u)-V d \bigr| \;=\; o(|d|)
  \qquad (d\to0),
\]
y \emph{fuertemente semisuave} si el miembro izquierdo es $O(|d|^2)$.
\end{definition}

\begin{lemma}
\label{lem:pl-semisuave}
Toda función lineal a trozos es fuertemente semisuave. Más aún, para $|d|$
suficientemente pequeño el término de error es exactamente $0$ si no hay quiebres
en el segmento $[u,u+d]$.
\end{lemma}
\begin{proof}
Fijado $u$, elíjase $r>0$ de modo que no haya quiebres distintos de $u$
en $(u-r,u+r)$. Para $0<|d|<r$, el punto $u+d$ pertenece al interior
de un tramo afín cuya clausura contiene a $u$. Por continuidad, si $V$
es su pendiente, $g(u+d)-g(u)=Vd$. En ese punto
$\dC g(u+d)=\{V\}$, por lo que el supremo de la definición es cero.
Para $d=0$ la igualdad es inmediata. Se obtiene así semisuavidad fuerte.
\end{proof}

\begin{theorem}[Convergencia del método de Newton semisuave, \citealp{qi1993}]
\label{thm:qi-sun}
Sea $g$ semisuave en una raíz $u^*$ y $BD$-regular allí, es decir, todo
$V\in\dB g(u^*)$ es no singular. Entonces existe $\varepsilon>0$ tal que la
iteración
\[
  u_{k+1}\;=\;u_k - V_k^{-1}g(u_k),\qquad V_k\in\dB g(u_k),
\]
está bien definida y converge superlinealmente a $u^*$ para todo
$u_0$ con $|u_0-u^*|<\varepsilon$, o termina en un número finito de pasos.
Si $g$ es fuertemente semisuave, satisface localmente
$|u_{k+1}-u^*|=O(|u_k-u^*|^2)$; esta cota incluye la terminación finita.
\end{theorem}

\section{Convergencia de métodos iterativos}
\label{sec:convergencia}

\begin{definition}[Orden de convergencia]
Una sucesión $\{x_k\}$ con $x_k\to x^*$ y $e_k=|x_k-x^*|>0$ converge con
\emph{orden} $p\ge1$ y \emph{constante asintótica} $C>0$ si
\[
  \lim_{k\to\infty}\frac{e_{k+1}}{e_k^{\,p}}\;=\;C.
\]
Para $p=1$ y $0<C<1$ se habla de convergencia \emph{lineal} de
razón $C$; para $p=2$ y $0<C<\infty$, de convergencia \emph{cuadrática}.
La convergencia es \emph{superlineal} si $e_{k+1}/e_k\to0$.
La terminación finita no se interpreta mediante un cociente asintótico con
errores estrictamente positivos.
\end{definition}

\begin{remark}[Descriptores empíricos de la convergencia]
\label{rem:estimacion-orden}
El estimador clásico \citep[cap.~2]{burden2016},
$\hat p_k=\ln(e_{k+1}/e_k)\big/\ln(e_k/e_{k-1})$, es muy sensible al ruido: basta
que un iterado caiga fortuitamente cerca de la raíz para que devuelva un valor sin
sentido. Esto ocurre sistemáticamente con la bisección, cuya sucesión de errores
\emph{no es monótona}. En este trabajo se calcula una pendiente descriptiva
por regresión de mínimos cuadrados de $\ln e_{k+1}$ contra $\ln e_k$ sobre
los pares de errores útiles, descartando los que están por debajo del ruido
de redondeo. Esa selección no identifica por sí sola una cola asintótica
ni certifica el orden teórico, especialmente si hay terminación en pocos
pasos. También se describen las razones entre errores sucesivos y, en los
métodos de encaje, la reducción del ancho del intervalo, que es la magnitud
que la teoría acota.
\end{remark}

\begin{definition}[Índice de eficiencia de Ostrowski, \citealp{ostrowski1966}]
\label{def:ostrowski}
Si un método tiene orden $p$ y requiere $c$ evaluaciones de la función objetivo
por iteración, su \emph{índice de eficiencia} es
\[
  E \;=\; p^{1/c}.
\]
\end{definition}

El índice de Ostrowski compara órdenes asintóticos cuando se realizan sus
hipótesis y el costo por evaluación es comparable. Con los órdenes clásicos
de la secante y de Newton con diferencias centradas se obtiene, respectivamente,
$E=\varphi\approx1{,}618$ y $E=2^{1/3}\approx1{,}260$. En este problema afín a
trozos, la terminación en pocos pasos impide interpretar esos valores como
eficiencias asintóticas observadas. La comparación empírica del \cref{cap:resultados} contabiliza las evaluaciones efectivamente realizadas.

\begin{theorem}[Cota de la bisección]
\label{thm:cota-biseccion}
Si $g$ es continua en $[a,b]$ y $g(a)g(b)<0$, la sucesión de puntos medios
$\{x_k\}$ de la bisección satisface
\[
  |x_k-x^*|\;\le\;\frac{b-a}{2^{k+1}},
\]
y por tanto alcanza una tolerancia $\varepsilon$ en a lo sumo
$\lceil\log_2((b-a)/\varepsilon)\rceil-1$ iteraciones.
\end{theorem}

La importancia del \cref{thm:cota-biseccion} para este trabajo es que su
demostración usa \emph{únicamente} la continuidad. Esta cota no requiere diferenciabilidad de $\VP$; por eso la bisección
funciona como referencia de robustez y como salvaguarda de los métodos híbridos.\footnote{Los métodos con encaje y las salvaguardas apropiadas también permiten garantías sin diferenciabilidad; debe distinguirse su cota de progreso de la cota específica de bisección.}
